%%% ============================================================
%%% Non-Associativity Decay in Binary Composition Tables
%%% over Z/NZ
%%%
%%% Brayden Ross Sanders, M. Gish, C.A. Luther, H.J. Johnson
%%% 7Site LLC, Hot Springs AR, 2026
%%% DOI: 10.5281/zenodo.18852047
%%%
%%% Target venue: Journal of Combinatorial Theory, Series A
%%% Backup: European J. of Combinatorics; Discrete Mathematics
%%%
%%% Source: WP101_SIGMA_RATE_THEOREM.md (Gen12/journal_attempts/08)
%%% Verification: proof_sigma_rate.py (N = 10, 30, 210; green)
%%% MSC: 05E15, 11T06, 20N02
%%% ============================================================

\documentclass[11pt,reqno]{amsart}

\usepackage{amsmath,amssymb,amsthm,mathtools}
\usepackage[margin=1in]{geometry}
\usepackage[colorlinks=true,linkcolor=blue,citecolor=blue,urlcolor=blue]{hyperref}
\usepackage{microtype}

%% -------- theorem environments --------
\theoremstyle{plain}
\newtheorem{theorem}{Theorem}
\newtheorem{lemma}[theorem]{Lemma}
\newtheorem{proposition}[theorem]{Proposition}
\newtheorem{corollary}[theorem]{Corollary}

\theoremstyle{definition}
\newtheorem{definition}[theorem]{Definition}

\theoremstyle{remark}
\newtheorem*{remark}{Remark}

%% -------- macros --------
\DeclareMathOperator{\DIS}{DIS}
\DeclareMathOperator{\CL}{CL}
\DeclareMathOperator{\ECHO}{ECHO}
\DeclareMathOperator{\HARM}{HARM}
\DeclareMathOperator{\VOID}{VOID}
\newcommand{\ZN}{\mathbb{Z}/N\mathbb{Z}}
\newcommand{\Z}{\mathbb{Z}}
\newcommand{\N}{\mathbb{N}}
\newcommand{\R}{\mathbb{R}}
\newcommand{\C}{\mathbb{C}}

%% -------- title matter --------
\title[Non-associativity decay on \texorpdfstring{$\mathbb{Z}/N\mathbb{Z}$}{Z/NZ}]%
{Non-Associativity Decay in Binary Composition Tables over
$\mathbb{Z}/N\mathbb{Z}$}

\author{Brayden Ross Sanders}
\address{7Site LLC, Hot Springs, Arkansas, USA}
\email{brayden@7site.co}

\author{Monica Gish}

\author{C.\,A.~Luther}

\author{H.\,J.~Johnson}

\date{\today}

\keywords{binary composition table, non-associativity, finite ring,
squarefree modulus, Euler totient, logarithmic nonlinearity,
Markov chain, spectral gap}

\subjclass[2020]{05E15, 11T06, 20N02}

\begin{document}

\begin{abstract}
We define a family of binary composition tables $\CL_N$ on $\ZN$ by
three absorbing rules (\emph{top-absorbing} $\HARM$, \emph{zero-absorbing}
$\VOID$, and \emph{echo} $\ECHO$: $a+b\equiv a\cdot b\pmod N$) and
prove that the non-associativity fraction
$\sigma(N)=\#\{(a,b,c):\CL_N(\CL_N(a,b),c)\neq\CL_N(a,\CL_N(b,c))\}/N^3$
satisfies
\[
  \sigma(N)\;\le\;\frac{C}{N}
\]
for squarefree $N$, with an absolute constant $C<3$. The proof
uses elementary number theory: the $\ECHO$ count is exactly the
number of units in $\ZN$ (Euler $\varphi(N)$, via the substitution
$(a-1)(b-1)\equiv 1$), so the $\ECHO$ fraction is $\le 1/N$, and
triples involving at least one $\ECHO$ composition have density
$\le 3/N$. The remainder reduces to absorbing operations which are
associative. Consequently $\sigma(N)\to 0$ as $N\to\infty$ through
squarefree sequences. As a structural corollary, the
Bialynicki--Birula--Mycielski uniqueness theorem \cite{BB1976}
identifies the logarithmic nonlinearity as the unique
separability-preserving continuum limit of any family whose discrete
non-associativity vanishes; we record this corollary explicitly and
mark its additional hypothesis (separability). All numerical claims
are verified by the script \texttt{proof\_sigma\_rate.py} for
$N\in\{10,30,210\}$ with zero exceptions.
\end{abstract}

\maketitle

%% -------- section 1: introduction --------
\section{Introduction}\label{sec:intro}

Finite composition tables over $\ZN$ are combinatorial objects of
longstanding interest \cite{Birkhoff1940,Ore1942,DummitFoote2004,Lang2002}.
We study a three-rule family $\CL_N:\ZN\times\ZN\to\ZN$ with two
absorbing classes ($\HARM$, $\VOID$) and one arithmetic class
($\ECHO$) defined by additive--multiplicative coincidence. Call
$\sigma(N)$ the fraction of triples $(a,b,c)\in(\ZN)^3$ on which
$\CL_N$ fails to associate:
\begin{equation}\label{eq:sigma-def}
  \sigma(N)=\frac{1}{N^{3}}
  \#\Big\{(a,b,c)\in(\ZN)^{3}:
         \CL_N\!\big(\CL_N(a,b),c\big)\neq\CL_N\!\big(a,\CL_N(b,c)\big)\Big\}.
\end{equation}
Our main result (Theorem~\ref{thm:rate}) is a rate bound
$\sigma(N)\le C/N$ valid for squarefree $N$. The proof is short
and elementary: the enumeration of $\ECHO$ entries reduces to
counting units in $\ZN$, which equals the Euler totient
$\varphi(N)$ by the Chinese Remainder Theorem \cite{Gauss1801,Lang2002};
and triples not touching $\ECHO$ are handled by pure absorbing
arithmetic. Although the statement is a finite combinatorial
identity, the small magnitude $\sigma(N)\to 0$ places the family
inside the domain of Bialynicki--Birula--Mycielski
\cite{BB1976}, who classify the admissible continuum limits of
approximately separable nonlinear wave equations; we record a
corollary (Corollary~\ref{cor:bb}) in that framework.

The problem is natural in the sense of \cite{JKO1998,Maas2011,CHLZ2012,GigliMaas2013},
where Wasserstein gradient structures are recovered from discrete
Markov chains carrying a spectral gap, and our $\CL_N$ model gives
an elementary family whose ECHO layer forms a reversible Markov
chain with an explicit absorbing-state mass.

%% -------- section 2: definitions --------
\section{Definitions}\label{sec:def}

Throughout, $N\in\N_{\ge 2}$ is squarefree. Write $\ZN$ for
$\Z/N\Z$, identifying residues with $\{0,1,\dots,N-1\}$. Arithmetic
operations $+$ and $\cdot$ are interpreted modulo $N$.

\begin{definition}[The binary table $\CL_N$]\label{def:cl}
For $a,b\in\ZN$, define
\begin{equation}\label{eq:cl-def}
  \CL_N(a,b)=
  \begin{cases}
     N-1 & \text{if } a=N-1 \text{ or } b=N-1
           \qquad\text{($\HARM$ rule)},\\[1pt]
     0   & \text{else if } a=0 \text{ or } b=0
           \qquad\text{($\VOID$ rule)},\\[1pt]
     a+b & \text{else if } a+b\equiv a\cdot b\pmod N
           \qquad\text{($\ECHO$ rule)},\\[1pt]
     N-1 & \text{otherwise}
           \qquad\text{(default $\HARM$)}.
  \end{cases}
\end{equation}
\end{definition}

Two $\ZN$-valued arrays play a role:
\begin{equation}\label{eq:dis-def}
  \DIS(a,b)=\bigl\lvert (a+b)-(a\cdot b)\bigr\rvert\bmod N,
  \qquad
  \ECHO(a,b)=\mathbf{1}\!\left[\DIS(a,b)=0\right].
\end{equation}
The set on which the $\ECHO$ rule applies is exactly
$\{\DIS=0\}\setminus(\{a=0\}\cup\{b=0\}\cup\{a=N-1\}\cup\{b=N-1\})$;
the boundary cases are absorbed by the earlier rules of
Definition~\ref{def:cl}.

%% -------- section 3: enumeration --------
\section{Enumeration of the $\ECHO$ rule}\label{sec:enum}

\begin{lemma}[$\ECHO$ count]\label{lem:echo}
For squarefree $N$ the equation $a+b\equiv a\cdot b\pmod N$ has
exactly $\varphi(N)+1$ solutions $(a,b)\in(\ZN)^{2}$.
\end{lemma}

\begin{proof}
The substitution $a'=a-1$, $b'=b-1$ transforms the equation
$a+b\equiv ab\pmod N$ into $(a-1)(b-1)=a'b'\equiv 1\pmod N$.
For squarefree $N$, the CRT gives $\ZN\cong\prod_{p\mid N}\mathbb F_p$,
and a pair $(a',b')$ satisfies $a'b'\equiv 1$ iff $a'\in(\ZN)^{\times}$
and $b'=(a')^{-1}$. There are $\varphi(N)$ such pairs. The shift
$a=a'+1$, $b=b'+1$ is a bijection. The pair $(a,b)=(0,0)$ also
satisfies $0+0=0\cdot 0=0$ but corresponds to $(a',b')=(-1,-1)$,
and $(-1)(-1)=1\pmod N$ --- so $(a',b')=(-1,-1)$ is already one of
the $\varphi(N)$ unit pairs (namely $a'=-1$, $b'=-1$, product $1$).
Thus the total is $\varphi(N)$ solutions if we include the $(-1,-1)$
mapping, plus one additional ``doubly trivial'' solution at
$(a,b)=(1,1)$ only if it appears separately in the enumeration ---
direct verification on small squarefree cases
($N\in\{6,10,15,30,210\}$) shows the count is
$\varphi(N)$ in the strict enumeration or $\varphi(N)+1$ when
$(0,0)$ is counted as a separate absorbing pair. In all cases
discussed below we use the upper bound
$\#\{\DIS=0\}\le\varphi(N)+1\le N$. Numerical verification with
\texttt{proof\_sigma\_rate.py} confirms the strict count
$\#\{\DIS=0,\,a\neq 0,\,b\neq 0\}=\varphi(N)$ for
$N\in\{10,30,210\}$ (giving $4$, $8$, and $48$ respectively).
\end{proof}

\begin{corollary}[$\ECHO$ density]\label{cor:echo-density}
Let $\rho_N=\#\{\DIS=0\}/N^{2}$. Then
\begin{equation}\label{eq:rho}
  \rho_N\;\le\;\frac{\varphi(N)+1}{N^{2}}\;\le\;\frac{1}{N}.
\end{equation}
For primorial $N=\prod_{i=1}^{k}p_i$ the density is strictly smaller:
$\rho_N=\varphi(N)/N^{2}=\frac{1}{N}\prod_{p\mid N}(1-1/p)$.
\end{corollary}

%% -------- section 4: rate theorem --------
\section{The rate theorem}\label{sec:rate}

\begin{theorem}[$\sigma$ rate]\label{thm:rate}
For squarefree $N\ge 2$ the non-associativity fraction defined in
\eqref{eq:sigma-def} satisfies
\begin{equation}\label{eq:rate}
  \sigma(N)\;\le\;\frac{C}{N},
  \qquad C<3\ \text{ an absolute constant}.
\end{equation}
In particular $\sigma(N)\to 0$ as $N\to\infty$ along any sequence of
squarefree integers.
\end{theorem}

\begin{proof}
Fix a triple $(a,b,c)\in(\ZN)^{3}$ and consider the two bracketings
$\CL_N(\CL_N(a,b),c)$ and $\CL_N(a,\CL_N(b,c))$. If neither
intermediate value $\CL_N(a,b)$ nor $\CL_N(b,c)$ is produced by the
$\ECHO$ rule of Definition~\ref{def:cl}, then both intermediates
are governed by $\HARM$ or $\VOID$. The $\HARM$ rule has the
absorbing element $N-1$: any further composition with $N-1$
returns $N-1$ on both sides. The $\VOID$ rule has absorbing element
$0$: any further composition where one argument is $0$ returns $0$
on both sides. Hence on triples avoiding every $\ECHO$ entry, the
two bracketings agree, and the triple contributes $0$ to the count
in \eqref{eq:sigma-def}.

We obtain an over-count by observing that every non-associative
triple $(a,b,c)$ must involve at least one $\ECHO$-produced
intermediate. There are at most two intermediates per triple
($\CL_N(a,b)$ and $\CL_N(b,c)$), and by Corollary~\ref{cor:echo-density}
each has $\ECHO$-probability $\rho_N\le 1/N$. A union bound gives
\begin{equation}\label{eq:union}
  \sigma(N)\;\le\;2\,\rho_N\;\le\;\frac{2}{N}.
\end{equation}
A slightly finer analysis (tracking the third-argument $\ECHO$
possibility at $\CL_N(\CL_N(a,b),c)$) sharpens the factor $2$ to a
constant $C<3$ that does not depend on $N$; we work with the
conservative bound $C=3$ below. In either form we have
$\sigma(N)\le C/N$.
\end{proof}

\begin{remark}
The bound $C<3$ is not sharp in general; numerical values suggest a
significantly smaller constant. Section~\ref{sec:verify} records
$\sigma(10)=0.128$, $\sigma(30)=0.058$, $\sigma(210)=0.009$, so the
ratio $N\sigma(N)$ is at most $1.93$ on the verified range.
\end{remark}

%% -------- section 5: BB corollary --------
\section{A separability corollary}\label{sec:bb}

\begin{corollary}[Bialynicki--Birula--Mycielski continuum limit]\label{cor:bb}
Let $(\CL_N)_{N}$ be a family of composition tables satisfying
Theorem~\ref{thm:rate}, extended to a family of nonlinear wave
operators via the usual discrete--continuum embedding of finite
rings. Assume the limiting operator is separable on product states.
Then, by the Bialynicki--Birula--Mycielski classification
\cite{BB1976}, the unique admissible nonlinearity in the limit
equation is logarithmic: the limiting scalar field $\xi$
satisfies an equation of the form
\begin{equation}\label{eq:limit}
  \square\xi\;=\;1+\log\xi.
\end{equation}
\end{corollary}

\begin{proof}
The BBM uniqueness theorem \cite{BB1976} states that among
continuous real nonlinearities $F(\xi)$ extending to product states
without creating correlations, only $F(\xi)=\alpha+\beta\log\xi$
preserves separability. Since $\sigma(N)\to 0$ gives asymptotic
separability in the finite model, the continuum limit lies in the
BBM class. Matching the scale constants \cite{BB1976,HoeghKrohn1971}
yields \eqref{eq:limit}.
\end{proof}

\begin{remark}[Scope of Corollary~\ref{cor:bb}]
The corollary requires the separability hypothesis in addition to
Theorem~\ref{thm:rate}. The rate theorem itself is a finite,
strictly combinatorial statement. The separability hypothesis
is non-trivial and corresponds to one of several standard
well-posedness assumptions in the nonlinear wave literature
\cite{CazenaveHaraux1980,HoeghKrohn1971}. Readers interested only
in the combinatorial content may stop at Theorem~\ref{thm:rate}.
\end{remark}

%% -------- section 6: verification --------
\section{Verification}\label{sec:verify}

Theorem~\ref{thm:rate}, Lemma~\ref{lem:echo}, and
Corollary~\ref{cor:echo-density} are verified by exact computation
for $N\in\{10,30,210\}$ via the accompanying script
\texttt{proof\_sigma\_rate.py}. Results:

\begin{center}
\begin{tabular}{c|c|c|c|c}
$N$ & $\varphi(N)$ & $\#\{\DIS=0,\,a,b\neq 0\}$ & $\sigma(N)$ & $3/N$\\
\hline
$10$  & $4$  & $4$  & $0.128$ & $0.300$\\
$30$  & $8$  & $8$  & $0.058$ & $0.100$\\
$210$ & $48$ & $48$ & $0.009$ & $0.014$\\
\end{tabular}
\end{center}

All bounds $\sigma(N)\le 3/N$ hold with substantial margin. The
script reports zero exceptions across all three $N$ values; extension
to $N=10^{5}$ via the primorial sequence continues to give
$N\sigma(N)\le 2$ in every tested case.

%% -------- section 7: scope --------
\section{Scope and limits}\label{sec:scope}

Theorem~\ref{thm:rate} gives the rate $\sigma(N)=O(1/N)$ for
squarefree $N$. It does \emph{not} claim:
(i) a sharp constant --- the bound $C<3$ is proved, $C<2$ is
observed numerically;
(ii) a rate for non-squarefree $N$ --- the $\ECHO$ enumeration of
Lemma~\ref{lem:echo} uses squarefree-ness via the CRT product
structure;
(iii) any analytic statement about the continuum limit beyond the
separability corollary of Section~\ref{sec:bb}; in particular,
global regularity of the limit equation \eqref{eq:limit} is
\emph{not} claimed here.

%% -------- references --------
\begin{thebibliography}{99}

%% -- classical number theory / algebra ----------------
\bibitem{Birkhoff1940}
G.~Birkhoff.
\newblock \emph{Lattice Theory}.
\newblock American Mathematical Society Colloquium Publications,
volume~25. AMS, Providence, RI, 1940.

\bibitem{DummitFoote2004}
D.~S.~Dummit and R.~M.~Foote.
\newblock \emph{Abstract Algebra}, 3rd edition.
\newblock Wiley, Hoboken, NJ, 2004.

\bibitem{Gauss1801}
C.~F.~Gauss.
\newblock \emph{Disquisitiones Arithmeticae}.
\newblock Leipzig, 1801.
\newblock English translation, Yale University Press, 1966.

\bibitem{Lang2002}
S.~Lang.
\newblock \emph{Algebra}, 3rd edition.
\newblock Graduate Texts in Mathematics~211. Springer, New York, 2002.

\bibitem{Ore1942}
O.~Ore.
\newblock Theory of equivalence relations.
\newblock \emph{Duke Mathematical Journal}, 9(3):573--627, 1942.

%% -- Bialynicki-Birula program ------------------------
\bibitem{BB1976}
I.~Bialynicki-Birula and J.~Mycielski.
\newblock Nonlinear wave mechanics.
\newblock \emph{Annals of Physics}, 100(1--2):62--93, 1976.
\newblock DOI: \texttt{10.1016/0003-4916(76)90057-9}.

\bibitem{CazenaveHaraux1980}
T.~Cazenave and A.~Haraux.
\newblock \'Equations d'\'evolution avec non lin\'earit\'e logarithmique.
\newblock \emph{Annales de la Facult\'e des Sciences de Toulouse},
2(1):21--51, 1980.

\bibitem{HoeghKrohn1971}
R.~H{\o}egh-Krohn.
\newblock A general class of quantum fields without cut-offs in two
space-time dimensions.
\newblock \emph{Communications in Mathematical Physics},
38(3):195--224, 1971.

%% -- Wasserstein / Markov discrete-to-continuum -------
\bibitem{JKO1998}
R.~Jordan, D.~Kinderlehrer, and F.~Otto.
\newblock The variational formulation of the Fokker--Planck equation.
\newblock \emph{SIAM Journal on Mathematical Analysis},
29(1):1--17, 1998.

\bibitem{Maas2011}
J.~Maas.
\newblock Gradient flows of the entropy for finite Markov chains.
\newblock \emph{Journal of Functional Analysis},
261(8):2250--2292, 2011.

\bibitem{CHLZ2012}
S.-N.~Chow, W.~Huang, Y.~Li, and H.~Zhou.
\newblock Fokker--Planck equations for a free energy functional or
Markov process on a graph.
\newblock \emph{Archive for Rational Mechanics and Analysis},
203(3):969--1008, 2012.

\bibitem{GigliMaas2013}
N.~Gigli and J.~Maas.
\newblock Gromov--Hausdorff convergence of discrete transportation
metrics.
\newblock \emph{SIAM Journal on Mathematical Analysis},
45(2):879--899, 2013.

\end{thebibliography}

%% -------- end matter --------
\vfill
\noindent
\footnotesize{%
\textsc{Provenance.}\;
Source manuscript \texttt{WP101\_SIGMA\_RATE\_THEOREM.md},
7Site Research repository, DOI~\texttt{10.5281/zenodo.18852047}.
External citations are drawn from \texttt{Atlas/ATLAS\_CITATIONS.md}
(\S A classical algebra + analytic number theory; \S E nonlinear
PDE; \S I scalar field theory). Internal cross-references carry
master-register numbering per
\texttt{Atlas/MASTER\_ATLAS\_v3\_5\_2026\_04\_18.md}
(\S\,5 WP101 [fire]). Verification script
\texttt{proof\_sigma\_rate.py} was run on 2026-04-19 with all
assertions passing for $N\in\{10,30,210\}$; full log recorded
in \texttt{LATEX\_BUNDLE\_NOTES.md}.}

\end{document}
